\section{Fazit}
\label{sec:conclusion}

Dieses Projekt hat untersucht, ob Hessian-basierte Krümmungsmetriken vorhersagen koennen, welche Schichten eines CNN unter der Power-of-Two-Quantisierung am stärksten geschädigt werden, und liefert auf alle drei Forschungsfragen eine klare, wenn auch nuancierte Antwort.
\par\medskip

Die Random-Init-Kontrolle (\S\ref{sec:phase2}) zeigt, dass die für die Analyse zentrale \texttt{conv1}-Kruemmungsspitze in allen drei Architekturen durch Training entsteht und nicht bereits bei Initialisierung vorhanden ist. Die Fusions-Analyse (\S\ref{sec:phase35}) zeigt weiter, dass diese Spitze überwiegend ein Effekt der Batch-Normalisierungs-Fusion ist und nicht der PoT-Quantisierung selbst, mit Ausnahme des \texttt{cnn}, wo PTQ zusätzlich moderat verstärkt. Die relative Krümmungsstruktur selbst bleibt zwischen PTQ- und QAT-Checkpoints nahezu unverändert ($\rho=0.95$--$1.00$ über alle sechs Kombinationen), waehrend die Beziehung dieser Krümmung zum tatsächlichen Schaden unter QAT fast vollständig zusammenbricht (0 von 18 signifikant, gegenüber 5 von 18 unter PTQ). Die isolierte Schichtquantisierung verbessert unter QAT häufig den Loss gegenüber der FP32-Baseline, statt ihn zu verschlechtern (\S\ref{sec:phase1-results}), sodass dieselbe Krümmungslandschaft ein grundlegend anderes Zielsignal vorherzusagen versucht (\S\ref{sec:disc-ptqqat}).
\par\medskip

Weiter ist kein einzelner Score universell zuverlässig, um die am stärksten geschädigte Schicht zu identifizieren, und selbst die rohe Krümmung ($S_{\mathrm{raw}}$) hält bei näherer Prüfung nur ein eng begrenztes Signal. Sie identifiziert zwar die wahre Top-Damage-Schicht in 3 von 6 Kombinationen (Rang~1), doch schließen 5 von 6 Bootstrap-Konfidenzintervallen Null ein, und im einzigen signifikanten Fall (ResNet-50/CIFAR10) liegt dieselbe Top-Schicht auf Rang 40 von 54 (Tabelle~\ref{tab:rank}). Eine Kontrollrechnung ohne \texttt{conv1} (Tabelle~\ref{tab:rank-noconv1}) zeigt, dass zwei der drei Rang-1-Erfolge direkt an \texttt{conv1} als wahre Top-Damage-Schicht gebunden sind, wobei ihr Rang nach dessen Entfernung auf 19 von 20 bzw.\ 50 von 53 faellt und dass \texttt{conv1}s Krümmungsspitze selbst überwiegend ein Fusions-/Fan-in-Artefakt ist (\S\ref{sec:phase35}). Nur der ResNet-18/CIFAR10-Erfolg übersteht die Kontrolle, weil seine Top-Damage-Schicht ohnehin nie \texttt{conv1} war. Ein Vergleich mit zwei Baselines ohne jede Hessian-Berechnung zeigt zudem, dass beide in allen sechs Kombinationen dieselbe oder eine leicht bessere Rang- und top-$k^*$-Trefferquote erreichen wie $S_{\mathrm{raw}}$, bei einem Bruchteil des Rechenaufwands. Das HAWQ-V2-Produkt scheitert zusaetzlich spezifisch daran, dass die PoT-Stoerungsgroeße mit der wahren Sensitivität anti-korreliert, statt neutral oder positiv zu sein (\S\ref{sec:disc-hawq}), was ein PoT-spezifischer Mechanismus und keine allgemeine Schwäche Krümmungs-basierter Sensitivität ist. Welche Schicht konkret am stärksten geschädigt wird, hängt zudem stark von Architektur und Datensatz ab (\texttt{conv1} auf ImageNet100 in allen drei Architekturen, aber uneinheitlich auf CIFAR10, Tabelle~\ref{tab:rank}).
\par\medskip

Letztens dominiert der Gewichtsschaden unter PTQ auf CIFAR10 in allen drei Architekturen, unter QAT dominiert der Aktivierungsschaden (\S\ref{sec:phase1-results}). Die Hessian-basierte Sensitivitätsanalyse, die ausschließlich Gewichtsschaden erfasst, ist damit nicht universell anwendbar, sondern regimeabhängig, und ihre Anwendbarkeit sollte vor jedem Einsatz explizit geprüft werden.
\par\medskip